\emph{Calepin} renders LaTeX math notation using KaTeX (HTML, optionally MathJax) or natively (LaTeX and Typst). Use \texttt{\textbackslash{}\$} for a literal dollar sign.

\subsection*{Inline and display math}
\label{inline-and-display-math}

Inline math is enclosed in single dollar signs: $a^2 + b^2 = c^2$.

Display math uses double dollar signs:

$$
\int_0^1 x^2 \, dx = \frac{1}{3}
$$

Multi-line equations use \texttt{\textbackslash{}begin\{aligned\}} inside double dollar signs:

\begin{equation}
\begin{aligned}
\nabla \cdot \mathbf{E} &= \frac{\rho}{\varepsilon_0} \\
\nabla \cdot \mathbf{B} &= 0 \\
\nabla \times \mathbf{E} &= -\frac{\partial \mathbf{B}}{\partial t} \\
\nabla \times \mathbf{B} &= \mu_0 \mathbf{J} + \mu_0 \varepsilon_0 \frac{\partial \mathbf{E}}{\partial t}
\end{aligned}
\label{eq-hello}
\end{equation}

And we can use standard cross-references as in \hyperref[eq-hello]{Equation 1}.

\subsection*{Theorem environments}
\label{theorem-environments}

Fenced divs with theorem-type classes get automatic numbering and can be cross-referenced.



\par\medskip\noindent\textbf{Theorem 1.}\label{thm-pythagoras} \itshape{} In a right triangle, the square of the hypotenuse equals the sum of the squares of the other two sides: $a^2 + b^2 = c^2$.

\upshape\medskip

\par\medskip\noindent\textit{ Proof.} Let a right triangle have legs $a$, $b$ and hypotenuse $c$. By constructing squares on each side and comparing areas, we obtain $a^2 + b^2 = c^2$.

 \hfill$\square$\medskip

By \hyperref[thm-pythagoras]{Theorem 1}, the relationship is fundamental to Euclidean geometry.

\subsubsection*{Supported types}
\label{supported-types}

The following theorem types are available, each with its own counter:

\texttt{theorem}, \texttt{lemma}, \texttt{corollary}, \texttt{proposition}, \texttt{conjecture}, \texttt{definition}, \texttt{example}, \texttt{exercise}, \texttt{solution}, \texttt{remark}, and \texttt{algorithm}.

\subsubsection*{Cross-references}
\label{cross-references}

Fenced divs with reserved id prefixes participate in the cross-reference system:

\medskip
\begin{tblr}{colspec={Q[l]Q[l]X[l]}, hline{1,2,Z}={solid}}
Class & Prefix & Example \\
\texttt{theorem} & \texttt{thm} & \texttt{::: \{.theorem \#thm-pythagoras\}} \\
\texttt{lemma} & \texttt{lem} & \texttt{::: \{.lemma \#lem-helper\}} \\
\texttt{definition} & \texttt{def} & \texttt{::: \{.definition \#def-group\}} \\
\texttt{example} & \texttt{exm} & \texttt{::: \{.example \#exm-basic\}} \\
\end{tblr}
\medskip

References like \texttt{\hyperref[thm-pythagoras]{Theorem 1}} resolve to the theorem number. The \texttt{\{\{label\}\}} variable generates format-specific anchors: \texttt{\textbackslash{}label\{thm-pythagoras\}} (LaTeX), \texttt{\textless{}thm-pythagoras\textgreater{}} (Typst).

